% \chapter{Atenção ao Estudante}

% O curso de Engenharia de Computação da UFU compreende que a excelência acadêmica exige um ambiente de aprendizagem acolhedor, saudável, diverso e acessível. Por isso, desenvolve e apoia ações de atenção estudantil em articulação com a FEELT e os órgãos centrais da universidade.

% \section{Acompanhamento Acadêmico e Pedagógico}

% A Coordenação do Curso e a Direção da FEELT acompanham ativamente o percurso dos(as) estudantes, prestando orientação individual, apoio à matrícula, aconselhamento pedagógico e mediação de conflitos acadêmicos. Também são promovidas ações de recepção de calouros e oficinas de ambientação universitária.

% \section{Assistência Estudantil e Políticas de Permanência}

% A Pró-Reitoria de Assistência Estudantil (PROAE) oferece suporte amplo aos(às) estudantes, com ações conduzidas por diretorias especializadas como a DIPAE (Diretoria de Inclusão, Promoção e Assistência Estudantil). Entre os programas e serviços, destacam-se:
% \begin{itemize}
%     \item Auxílios financeiros para moradia, alimentação, transporte e creche;
%     \item Bolsas acadêmicas (monitoria, PIBIC, extensão, PET);
%     \item Apoio à acessibilidade e à inclusão de pessoas com deficiência;
%     \item Acompanhamento psicopedagógico e social individualizado;
%     \item Incentivo à inclusão digital e ao acesso universal ao conhecimento;
%     \item Reconhecimento da diversidade étnico-racial, cultural e de gênero.
% \end{itemize}

% \section{Saúde Mental e Bem-Estar Psicológico}

% A UFU mantém serviços de orientação em saúde mental por meio da Divisão de Saúde do Estudante, que atua em parceria com o Sistema de Saúde Universitário (SIS-UFU) e projetos de promoção de bem-estar. O atendimento contempla:
% \begin{itemize}
%     \item Apoio psicológico individual e escuta qualificada;
%     \item Encaminhamento para rede SUS quando necessário;
%     \item Campanhas educativas e preventivas (ex: ansiedade, depressão, suicídio);
%     \item Oficinas temáticas e grupos de apoio organizados pelo setor especializado.
% \end{itemize}

% Mais informações: \url{https://proae.ufu.br/servicos/orientacao-em-saude-mental}.

% \section{Promoção da Diversidade, Inclusão e Direitos Humanos}

% A UFU reafirma seu compromisso com o combate a todas as formas de preconceito, discriminação e exclusão. Campanhas institucionais e programas permanentes promovem o respeito às diferenças e a construção de um ambiente universitário mais justo, seguro e plural.

% Ações voltadas ao combate à LGBTfobia, ao racismo, ao capacitismo e à xenofobia são desenvolvidas em articulação com o Comitê de Diversidade da UFU, os coletivos estudantis e os programas institucionais da PROAE.

% \section{Canais de Apoio ao Estudante}

% A Universidade Federal de Uberlândia oferece diversos serviços de apoio acadêmico, psicossocial e institucional para garantir a permanência qualificada dos(as) estudantes. Abaixo, estão relacionados alguns dos principais canais de apoio:

% \begin{longtblr}[
%     theme = ecp,
%     caption = {Carga Horária por Componente Curricular},
%     label = {tab:carga_horaria_por_componente_curricular},
% ]{
% colspec = {Q[l,m,wd=30mm]Q[c,m,wd=40mm]Q[c,m,wd=30mm]},
% rowhead = 1,
% row{odd} = {bg=CinzaClaro},
% row{1} = {bg=AzulEscuro, fg=white},
% row{7} = {bg=AzulClaro, fg=white},
% cells  = {font = \fontsize{10pt}{12pt}\selectfont},
% }
% \textbf{Área de Apoio} & \textbf{Serviço / Portal} & \textbf{Link}\\
% Apoio pedagógico individual e coletivo & Divisão de Promoção de Igualdades e Apoio Educacional (DIPAE) & \url{https://proae.ufu.br/dipae}\\
% Apoio psicossocial e saúde mental & Orientação e atendimento psicológico & proae.ufu.br/servicos/orientacao-em-saude-mental\\
% Inclusão, permanência e assistência estudantil & Diretoria de Inclusão, Promoção e Assistência Estudantil (DIPAE) & proae.ufu.br/dipae\\
% Ações de respeito e diversidade & Campanhas e ações institucionais & comunica.ufu.br/noticias/2025/05/respeito-e-diversidade\\
% Sistema de Saúde Universitário & Atendimento médico, odontológico e psicológico (SIS-UFU) & sis.ufu.br\\
% Mobilidade acadêmica nacional &  & \url{https://prograd.ufu.br/mobilidadenacional}\\
% Mobilidade acadêmica internacional & Diretoria de Relações Internacionais e Interinstitucionais (DRII) & \url{https://dri.ufu.br/}\\
% Sistemas e documentos acadêmicos & Portal do Estudante de Graduação & \url{https://www.portalestudante.ufu.br/}\\
% Suporte à vida acadêmica & Coordenação do Curso de \nomecursonovo (FEELT) & \url{https://www.feelt.ufu.br/graduacao/engenharia-de-computacao}\\
% Portal de notícias & Informações institucionais e notícias & \url{https://comunica.ufu.br/}\\
% Ouvidoria e denúncias & Canal de Ouvidoria & \url{https://ufu.br/ouvidoria-geral}\\
% \end{longtblr}

% \textbf{Observação:} A UFU também conta com comissões permanentes, grupos estudantis, projetos de extensão e programas de mentoria voluntária que fortalecem a escuta e o acolhimento à comunidade discente.

% \section{Canais de Apoio ao Estudante}

% A Universidade Federal de Uberlândia disponibiliza uma ampla rede de apoio aos(às) estudantes, com foco na permanência qualificada, saúde, inclusão e bem-estar. Para informações detalhadas, editais atualizados e inscrições, recomenda-se acompanhar regularmente:

% \begin{itemize}
%     \item O site da Pró-Reitoria de Assistência Estudantil (PROAE):
    
%     \url{https://proae.ufu.br/tags/auxilios-estudantis}.

%     \item Editais PROAE publicados periodicamente e que contêm informações sobre auxílios diversos:
    
%     \url{https://proae.ufu.br/tags/auxilios-estudantis}
    
%     \item Fórum de Assuntos Estudantis, um espaço permanente de debate, proposição, negociação, reivindicação e encaminhamento de demandas referentes a política de assistência estudantil da universidade:
    
%     \url{https://proae.ufu.br/servicos/forum-de-assuntos-estudantis-fae}
% \end{itemize}
    

%     \item Programa de Bolsa Permanência (vinculado ao MEC):
    
%     \url{https://proae.ufu.br/unidades-organizacionais/diretoria-de-inclusao-promocao-e-assistencia-estudantil}

% Assistência e Orientação Social (DIASE): \url{https://proae.ufu.br/diase}
% Apoio Educacional e Promoção de Igualdades (DIPAE): \url{https://proae.ufu.br/dipae}
% Esporte e Lazer (DIESU): \url{https://proae.ufu.br/diesu}
% Moradia Estudantil (DIVME): \url{https://proae.ufu.br/divme}
% Restaurantes universitários (RU): \url{https://proae.ufu.br/ru}
% Saúde do Estudante (DISAU): \url{https://proae.ufu.br/disau}


% \subsection{Auxílios Financeiros e de Permanência}
% \begin{longtblr}[
%     theme = ecp,
%     caption = {Carga Horária por Componente Curricular},
%     label = {tab:carga_horaria_por_componente_curricular},
% ]{
% colspec = {Q[l,m,wd=40mm]Q[c,m,wd=30mm]Q[c,m,wd=30mm]},
% rowhead = 1,
% row{odd} = {bg=CinzaClaro},
% row{1} = {bg=AzulEscuro, fg=white},
% row{7} = {bg=AzulClaro, fg=white},
% cells  = {font = \fontsize{10pt}{12pt}\selectfont},
% }
% Programa / Auxílio & Descrição & Base normativa\\
% Bolsa Permanência & Apoio financeiro contínuo para estudantes em vulnerabilidade socioeconômica, com vínculo ao MEC & Portaria MEC nº 22/2024\\
% Auxílio Moradia & Concessão mensal para estudantes que residem longe da unidade de ensino e não possuem condições de manter moradia própria & Editais PROAE\\
% Auxílio Alimentação & Apoio financeiro ou refeição subsidiada no Restaurante Universitário & Editais PROAE\\
% Auxílio Transporte & Subsídio para deslocamento de estudantes entre residência e universidade & Editais PROAE\\
% Auxílio Competição Esportiva & Apoio financeiro para participação em competições esportivas representando a UFU & Editais PROAE\\
% Auxílio Inclusão Digital & Auxílio temporário para aquisição de equipamentos ou acesso à internet & Programa PIID (vinculado ao PROMISAES)\\
% \end{longtblr}

% \subsection{Apoio Psicossocial e à Saúde Mental}
% \begin{longtblr}[
%     theme = ecp,
%     caption = {Carga Horária por Componente Curricular},
%     label = {tab:carga_horaria_por_componente_curricular},
% ]{
% colspec = {Q[l,m,wd=40mm]Q[c,m,wd=30mm]Q[c,m,wd=30mm]},
% rowhead = 1,
% row{odd} = {bg=CinzaClaro},
% row{1} = {bg=AzulEscuro, fg=white},
% row{7} = {bg=AzulClaro, fg=white},
% cells  = {font = \fontsize{10pt}{12pt}\selectfont},
% }
% Programa / Serviço & Descrição & Acesso\\
% Acompanhamento psicológico individual & Atendimento e orientação em saúde mental com psicólogos(as) da Divisão de Saúde do Estudante & proae.ufu.br/servicos/orientacao-em-saude-mental\\
% Encaminhamentos para o SUS & Articulação com a rede pública de saúde para casos que exigem acompanhamento contínuo & SIS-UFU\\
% Campanhas de prevenção & Ações temáticas sobre ansiedade, depressão, prevenção ao suicídio, entre outros & PROAE / DIPAE\\
% \end{longtblr}

% \subsection{Apoio à Diversidade, Inclusão e Internacionalização}
% \begin{longtblr}[
%     theme = ecp,
%     caption = {Carga Horária por Componente Curricular},
%     label = {tab:carga_horaria_por_componente_curricular},
% ]{
% colspec = {Q[l,m,wd=40mm]Q[c,m,wd=30mm]Q[c,m,wd=30mm]},
% rowhead = 1,
% row{odd} = {bg=CinzaClaro},
% row{1} = {bg=AzulEscuro, fg=white},
% row{7} = {bg=AzulClaro, fg=white},
% cells  = {font = \fontsize{10pt}{12pt}\selectfont},
% }
% Programa / Ação & Descrição & Base normativa\\
% PROMISAES (Programa de Acesso ao Ensino Superior) & Concessão de auxílio financeiro a estudantes estrangeiros em situação de vulnerabilidade & Editais PROAE\\
% Apoio a estudantes com deficiência (PcD) & Acompanhamento pedagógico e acessibilidade & DIPAE / NAI\\
% Campanhas contra discriminação & Ações institucionais de enfrentamento à LGBTfobia, racismo, xenofobia e capacitismo & Comissão Permanente de Acompanhamento da Política de Diversidade Sexual e de Gênero (CPDiversa)\\
% \end{longtblr}

% \subsection{Outros Programas Institucionais}
% \begin{longtblr}[
%     theme = ecp,
%     caption = {Carga Horária por Componente Curricular},
%     label = {tab:carga_horaria_por_componente_curricular},
% ]{
% colspec = {Q[l,m,wd=40mm]Q[c,m,wd=30mm]Q[c,m,wd=30mm]},
% rowhead = 1,
% row{odd} = {bg=CinzaClaro},
% row{1} = {bg=AzulEscuro, fg=white},
% row{7} = {bg=AzulClaro, fg=white},
% cells  = {font = \fontsize{10pt}{12pt}\selectfont},
% }
% Programa & Descrição & Acesso\\
% PET Engenharia de Computação & Formação tutorial integrada ao ensino, pesquisa e extensão & Coordenação PET\\
% Monitorias, Iniciação Científica e Extensão & Bolsas acadêmicas com inserção em projetos orientados por docentes & SIGProj / PIBIC / PIVIC\\
% Mobilidade acadêmica nacional e internacional & Programas como ANDIFES, Brafitec, e convênios internacionais da UFU & ari.ufu.br\\
% \end{longtblr}
    

\chapter{Atenção ao Estudante}

O curso de \nomecursonovo da UFU compreende que a excelência acadêmica exige um ambiente de aprendizagem acolhedor, saudável, diverso e acessível. Por isso, desenvolve e apoia ações de atenção estudantil em articulação com a Faculdade de Engenharia Elétrica (FEELT) e os órgãos centrais da universidade.

\section{Acompanhamento Acadêmico e Pedagógico}

A Coordenação do Curso e a Direção da FEELT acompanham ativamente o percurso dos(as) estudantes, prestando orientação individual, apoio à matrícula, aconselhamento pedagógico e mediação de conflitos acadêmicos. Também são promovidas ações de recepção de calouros e oficinas de ambientação universitária.

Para fortalecer esse suporte, a FEELT instituiu um núcleo específico para o acolhimento e a orientação da comunidade discente.

\section{Núcleo de Apoio e Atenção ao Estudante da FEELT (NAAE-FEELT)}

Em conformidade com a política institucional de assistência estudantil~\cite{UFU:CONSEX:20:2022} e seu regimento interno~\cite{UFU:FEELT:Regimento:2023}, a Faculdade de Engenharia Elétrica dispõe do Núcleo de Apoio e Atenção ao Estudante (NAAE-FEELT). Este núcleo atua como um ponto de referência para o acolhimento, a escuta qualificada e o encaminhamento das demandas estudantis no âmbito da faculdade.

Os principais objetivos do NAAE-FEELT são:
\begin{itemize}
\item Propor e executar ações de acolhimento e integração para estudantes ingressantes e veteranos;
\item Oferecer suporte pedagógico e psicossocial em primeira instância, atuando de forma preventiva e mediadora;
\item Articular as demandas dos estudantes com as coordenações de curso e os serviços especializados da UFU, como a Pró-Reitoria de Assistência Estudantil (PROAE);
\item Contribuir para a redução dos índices de retenção e evasão, promovendo a permanência qualificada;
\item Fomentar um ambiente acadêmico inclusivo, plural e respeitoso.
\end{itemize}

As ações a serem desenvolvidas incluem recepção de calouros, rodas de conversa, oficinas sobre métodos de estudo e saúde mental, além de atendimentos individuais e coletivos.

Mais informações:

{\footnotesize\url{https://www.feelt.ufu.br/unidades/nucleo/nucleo-de-apoio-e-atencao-ao-estudante}}

\section{Assistência Estudantil e Políticas de Permanência}

A Pró-Reitoria de Assistência Estudantil (PROAE) oferece suporte amplo aos(às) estudantes, com ações conduzidas por diretorias especializadas como a DIRES (Diretoria de Inclusão, Promoção e Assistência Estudantil) e DIPAE (Divisão de Promoção de Igualdade e Apoio Educacional). Entre os programas e serviços, destacam-se:

\begin{itemize}
\item Auxílios financeiros para moradia, alimentação, transporte e creche;
\item Bolsas acadêmicas (monitoria, PIBIC, extensão, PET);
\item Apoio à acessibilidade e à inclusão de pessoas com deficiência;
\item Acompanhamento psicopedagógico e social individualizado;
\item Incentivo à inclusão digital e ao acesso universal ao conhecimento;
\item Reconhecimento da diversidade étnico-racial, cultural e de gênero.
\end{itemize}

\section{Saúde Mental e Bem-Estar Psicológico}

A UFU mantém serviços de orientação em saúde mental por meio da Divisão de Saúde do Estudante (DISAU), que atua em parceria com o Sistema Único de Saúde (SUS) e projetos de promoção de bem-estar. O atendimento contempla:

\begin{itemize}
\item Apoio psicológico individual e escuta qualificada;
\item Encaminhamento para a rede SUS quando necessário;
\item Campanhas educativas e preventivas (ex: ansiedade, depressão, suicídio);
\item Oficinas temáticas e grupos de apoio organizados pelo setor especializado.
\end{itemize}

Mais informações:

{\footnotesize\url{https://proae.ufu.br/servicos/orientacao-em-saude-mental}}

\section{Promoção da Diversidade, Inclusão e Direitos Humanos}

A UFU reafirma seu compromisso com o combate a todas as formas de preconceito, discriminação e exclusão. Campanhas institucionais e programas permanentes promovem o respeito às diferenças e a construção de um ambiente universitário mais justo, seguro e plural.

Ações voltadas ao combate à LGBTfobia, ao racismo, ao capacitismo e à xenofobia são desenvolvidas em articulação com o Comitê de Diversidade da UFU, os coletivos estudantis e os programas institucionais da PROAE, DIEPAFRO e CPDiversa.

\section{Canais de Apoio ao Estudante}

A Universidade Federal de Uberlândia oferece diversos serviços de apoio acadêmico, psicossocial e institucional para garantir a permanência qualificada dos(as) estudantes. Abaixo, estão relacionados alguns dos principais canais de apoio:

\begin{longtblr}[
theme = ecp,
caption = {Principais Canais de Apoio ao Estudante},
label = {tab:canais_apoio_estudante},
]{
colspec = {Q[l,m,wd=30mm]Q[c,m,wd=35mm]Q[c,m,wd=70mm]},
rowhead = 1,
row{odd} = {bg=CinzaClaro},
row{1} = {bg=AzulEscuro, fg=white},
% row{7} = {bg=AzulClaro, fg=white},                                                             
cells  = {font = \fontsize{10pt}{12pt}\selectfont},
}
\textbf{Área de Apoio} & \textbf{Serviço / Portal} & \textbf{Link / Contato}\\
Apoio Local (FEELT) & Núcleo de Apoio e Atenção ao Estudante (NAAE-FEELT) & \url{https://www.feelt.ufu.br/unidades/nucleo/nucleo-de-apoio-e-atencao-ao-estudante}\\
Suporte à vida acadêmica & Coordenação do Curso de \nomecursonovo (FEELT) & \url{https://www.feelt.ufu.br/graduacao/engenharia-de-computacao}\\
Apoio pedagógico individual e coletivo & Divisão de Promoção de Igualdades e Apoio Educacional (DIPAE/PROAE) & \url{https://proae.ufu.br/dipae}\\
Apoio psicossocial e saúde mental & Divisão de Saúde do Estudante (DISAU/PROAE) & \url{https://proae.ufu.br/servicos/orientacao-em-saude-mental}\\
Inclusão, permanência e assistência estudantil & Diretoria de Inclusão, Promoção e Assistência Estudantil (DIRES) e Divisão de Promoção de Igualdade e Apoio Educacional (DIPAE) & \url{https://proae.ufu.br/dipae}\\
Ações de respeito e diversidade & DIEPAFRO, CPDiversa, Campanhas e ações institucionais & \url{https://comunica.ufu.br}\\
Sistema de Saúde Universitário & Divisão de Saúde do Estudante (DISAU/PROAE) & \url{https://proae.ufu.br/disau}\\
Mobilidade acadêmica nacional & Programa ANDIFES & \url{https://prograd.ufu.br/mobilidadenacional}\\
Mobilidade acadêmica internacional & Diretoria de Relações Internacionais e Interinstitucionais (DRII) & \url{https://dri.ufu.br/}\\
Sistemas e documentos acadêmicos & Portal do Estudante de Graduação & \url{https://www.portalestudante.ufu.br/}\\
Portal de notícias & Informações institucionais e notícias & \url{https://comunica.ufu.br/}\\
Ouvidoria e denúncias & Canal de Ouvidoria & \url{https://ufu.br/ouvidoria-geral}\\
\end{longtblr}

A UFU conta com comissões permanentes, grupos estudantis, projetos de extensão e programas de mentoria voluntária que fortalecem a escuta e o acolhimento à comunidade discente. É fundamental que o(a) estudante acompanhe os canais oficiais da PROAE e da FEELT para informações sobre editais e cronogramas.